% appendix/questions/after.tex
% mainfile: ../../perfbook.tex
% SPDX-License-Identifier: CC-BY-SA-3.0

\section{"之后"是什么意思？}
\label{sec:app:questions:What Does ``After'' Mean?}
%
\epigraph{没有未来。
	  没有过去。}
	 {Alan Moore}

"之后"是一个直观但却出人意料地困难的概念。
一个重要的非直观问题是，代码可以在任意时刻被延迟任意长的时间。
考虑一个生产者线程和一个消费者线程，它们通过一个全局结构体进行通信，
该结构体包含时间戳\qco{t}和整数字段\qco{a}、\qco{b}和\qco{c}。
生产者循环记录当前时间（以自1970年以来的秒数的十进制形式），
然后更新\qco{a}、\qco{b}和\qco{c}的值，
如\cref{lst:app:questions:After Producer Function}所示。
消费者代码也循环记录当前时间，
同时复制生产者的时间戳以及字段\qco{a}、\qco{b}和\qco{c}，
如\cref{lst:app:questions:After Consumer Function}所示。
在运行结束时，消费者输出异常记录的列表，
例如时间似乎倒流的情况。

\begin{listing}
\input{CodeSamples/api-pthreads/QAfter/time@producer.fcv}
\caption{"之后"生产者函数}
\label{lst:app:questions:After Producer Function}
\end{listing}

\begin{listing}
\ebresizeverb{.96}{\input{CodeSamples/api-pthreads/QAfter/time@consumer.fcv}}
\caption{"之后"消费者函数}
\label{lst:app:questions:After Consumer Function}
\end{listing}

\QuickQuiz{
	在这些示例中你能发现哪些SMP编程错误？
	完整代码请参见\path{time.c}。
}\QuickQuizAnswer{
	以下是你可能发现的错误：

	\begin{enumerate}
	\item	紧凑循环中缺少\co{barrier()}或\co{volatile}。
	\item	更新端缺少内存屏障。
	\item	生产者和消费者之间缺乏同步。
	\end{enumerate}
}\QuickQuizEnd

人们可能直觉地认为生产者和消费者时间戳之间的差异会非常小，
因为生产者记录时间戳或值不需要太长时间。
在双核1\,GHz x86上的一些样本输出摘录如
\cref{tab:app:questions:After Program Sample Output}所示。
其中，"seq"列是循环执行的次数，
"time"列是异常发生的时间（以秒为单位），
"delta"列是消费者的时间戳晚于生产者时间戳的秒数
（负值表示消费者在生产者之前收集了时间戳），
标记为"a"、"b"和"c"的列显示了这些变量
自消费者上次快照收集以来的增长量。

\begin{table}
\rowcolors{1}{}{lightgray}
\renewcommand*{\arraystretch}{1.2}
\sisetup{group-digits=false}
\centering
\scriptsize
\begin{tabular}{rS[table-format=7.6]rS[table-format=3.0]S[table-format=3.0]S[table-format=3.0]}
\toprule
seq    & \multicolumn{1}{c}{time (seconds)} & delta~  &  a &  b &  c \\
\midrule
17563: & 1152396.251585 & ($-16.928$) & 27 & 27 & 27 \\
18004: & 1152396.252581 & ($-12.875$) & 24 & 24 & 24 \\
18163: & 1152396.252955 & ($-19.073$) & 18 & 18 & 18 \\
18765: & 1152396.254449 & ($-148.773$) & 216 & 216 & 216 \\
19863: & 1152396.256960 & ($-6.914$) & 18 & 18 & 18 \\
21644: & 1152396.260959 & ($-5.960$) & 18 & 18 & 18 \\
23408: & 1152396.264957 & ($-20.027$) & 15 & 15 & 15 \\
\bottomrule
\end{tabular}
\caption{"之后"程序样本输出}
\label{tab:app:questions:After Program Sample Output}
\end{table}

为什么时间会倒流？
括号中的数字是以微秒为单位的差值，
其中一个大的数字超过了10微秒，甚至有一个超过了100微秒！
请注意，这个CPU在这段时间内可能执行超过100,000条指令。

\begin{fcvref}[ln:api-pthreads:QAfter:time]
一个可能的原因如下面的事件序列所示：
\begin{enumerate}
\item	消费者获取时间戳
	（\cref{lst:app:questions:After Consumer Function}，
	\clnref{consumer:tod}）。
\item	消费者被抢占。
\item	经过任意长的时间。
\item	生产者获取时间戳
	（\cref{lst:app:questions:After Producer Function}，
	\clnref{producer:tod}）。
\item	消费者重新开始运行，并获取生产者的时间戳
	（\cref{lst:app:questions:After Consumer Function}，
	\clnref{consumer:prodtod}）。
\end{enumerate}

在这种场景中，生产者的时间戳可能比消费者的时间戳晚任意长的时间。

如何在SMP代码中避免为"之后"的含义而烦恼？

只需按照设计使用SMP原语即可。

在这个示例中，最简单的修复方法是使用锁，
例如，在\cref{lst:app:questions:After Producer Function}中
\clnref{producer:tod}之前让生产者获取锁，
在\cref{lst:app:questions:After Consumer Function}中
\clnref{consumer:tod}之前让消费者获取锁。
该锁还必须在\cref{lst:app:questions:After Producer Function}中
\clnref{producer:upd:c}之后释放，
以及在\cref{lst:app:questions:After Consumer Function}中
\clnref{consumer:upd:c}之后释放。
这些锁使得\cref{lst:app:questions:After Producer Function}中
\clnrefrange{producer:tod}{producer:upd:c}的代码段和
\cref{lst:app:questions:After Consumer Function}中
\clnrefrange{consumer:tod}{consumer:upd:c}的代码段
相互{\em 排斥}，换句话说，相对于彼此原子地运行。
这在\cref{fig:app:questions:Effect of Locking on Snapshot Collection}
中有所表示：
锁防止了任何代码方框在时间上重叠，因此消费者的时间戳必须在
先前生产者的时间戳之后收集。
此图中每个方框中的代码段被称为"临界区"；
在给定时间只能有一个这样的临界区在执行。
\end{fcvref}

\begin{figure}
\centering
\includegraphics{appendix/questions/after-snapshot}
\caption{锁对快照收集的影响}
\label{fig:app:questions:Effect of Locking on Snapshot Collection}
\end{figure}

加入锁后，输出结果如
\cref{fig:app:questions:Locked After Program Sample Output}所示。
这里没有时间倒流的情况，
取而代之的是消费者连续读取之间计数差异超过1,000的情况。

\begin{table}
\renewcommand*{\arraystretch}{1.2}
\sisetup{group-digits=false}
\centering
\scriptsize
\begin{tabular}{rS[table-format=7.6]rS[table-format=4.0]S[table-format=4.0]S[table-format=4.0]}
\toprule
seq    & \multicolumn{1}{c}{time (seconds)} & delta~  &  a &  b &  c \\
\midrule
58597:  & 1156521.556296 & ($3.815$) & 1485 & 1485 & 1485 \\
403927: & 1156523.446636 & ($2.146$) & 2583 & 2583 & 2583 \\
\bottomrule
\end{tabular}
\caption{加锁后的"之后"程序样本输出}
\label{fig:app:questions:Locked After Program Sample Output}
\end{table}

\QuickQuiz{
	消费者连续读取之间怎么会有这么大的间隔？
	完整代码请参见\path{timelocked.c}。
}\QuickQuizAnswer{
	以下是出现这种间隔的一些原因：

	\begin{enumerate}
	\item	消费者可能被抢占很长时间。
	\item	一个长时间运行的中断可能延迟消费者。
	\item	缓存未命中可能延迟消费者。
	\item	生产者可能运行在比消费者更快的CPU上
		（例如，其中一个CPU可能由于散热或功耗约束
		不得不降低其时钟频率）。
	\end{enumerate}
}\QuickQuizEnd

总之，如果你获取了一个\IXh{exclusive}{lock}，你就{\em 知道}
在持有该锁期间你做的任何事情都将显示为发生在
该锁的任何先前持有者所做的事情之后，至少在
\IXacrl{tle}的误差范围内
（参见\cref{sec:future:Semantic Differences}）。
不需要担心哪个CPU是否执行了\IX{memory barrier}，
不需要担心CPU或编译器对操作的重排序——生活很简单。
当然，这种锁定阻止了这两段代码的并发运行，
可能限制了程序在多处理器上获得性能提升的能力，
可能导致"安全但缓慢"的情况。
\Cref{chp:Partitioning and Synchronization Design}描述了
在许多情况下获得性能和可扩展性的方法。

简言之，在许多并行程序中，"之后"的真正重要定义是操作的排序，
这在\cref{chp:Advanced Synchronization: Memory Ordering}
中有详尽的讨论。

然而，在大多数情况下，如果你发现自己在担心某段代码之前或之后会发生什么，
你应该把这看作一个提示，表明需要更好地使用标准原语。
让这些原语替你操心吧。
